\subsection{Parameters and statement}

Throughout this section assume
\begin{equation}\label{minimum:range}
 \frac18\le a\le\frac16,\qquad 2a\le b\le\frac12.
\end{equation}
Define
\begin{align*}
 k&=2b-a,&q&=\frac b{k},&h&=\frac{1+q}{2},\\
 \lambda&=\frac{1-2b}{4b},&L_0&=\frac{1-2a}{4},&
 E_1&=\frac{1-q}{3},\\
 E_2&=\frac{(b-a)(1+6b-4a)}{6k},&&&
 A_E&=\frac{1-a-3b}{3},\\
 \sigma_E&=\frac{6b-1-2a}{3a},&
 S_F&=\frac{1-a}{3}-h,&\sigma_F&=\frac{a+2}{3a},\\
 S_V&=\frac{1-a}{3}-\frac{(1-b)h}{2},&
 \sigma_V&=\frac{1-a}{6a}.&&
\end{align*}
All three quantities \(\sigma_E,\sigma_F,\sigma_V\) are strictly positive.
For \(\sigma_E\), use \(6b-1-2a\ge10a-1\ge1/4\).
Set
\begin{align}
 J_E&=\frac{\sigma_E L_0+\lambda A_E}{\sigma_E+\lambda}
     =\frac{(1-a)(b-a)(6b-1)}{24b^2-14ab+3a-4b},
       \label{minimum:JE}\\
 J_F&=\frac{\sigma_F L_0+\lambda S_F}{\sigma_F+\lambda},
       \label{minimum:JF}\\
 J_V&=\frac{\sigma_V L_0+\lambda S_V}{\sigma_V+\lambda}.
       \label{minimum:JV}
\end{align}
In particular, all the denominators in these expressions are positive.
The refined cost is the explicit continuous function
\begin{equation}\label{minimum:cost}
 V(a,b)=\max\bigl\{\min(E_2,J_E),\ \min(J_F,J_V)\bigr\}.
\end{equation}

\begin{lemma}[Profile bound retaining two actual minima]\label{minimum:profile}
Assume \eqref{minimum:range}. Let \(g:[0,N]\to\mathbb R\) be
1-Lipschitz, with \(ax\le g(x)\le bx\). If
\(7N/8\le n_0<N\), there is an admissible decreasing chain from
\(n_0\) to zero, without a prescribed intermediate endpoint, whose cost
is at most \(NV(a,b)\). Its number of edges satisfies
\[
 K\le C\left(1+\log\frac{N}{N-n_0}\right)
\]
with an absolute constant \(C\). On a grid of mesh \(T\) containing
\(0,n_0,N\), the additional cost is \(O(KT)\). Additive barrier errors
of size \(\epsilon N\) increase this to
\(O(K\epsilon N+KT)\).
\end{lemma}

The proof is elementary. It uses an affine potential, the location and value
of a minimum on \([q,1]\), and a reflection that raises successive new
minima to the same height. We provide the construction, including its
bounded-length and stability assertions.

\subsection{Two reusable chain constructions}
Normalize \(N=1\). Write \(\operatorname{Var}^-(g;I)\) for the negative
variation of \(g\) on an interval \(I\), in the direction of increasing
coordinate. Every edge cost is at most this negative variation on its
interval. Since a Lipschitz function is absolutely continuous,
\begin{equation}\label{minimum:variation}
 \operatorname{Var}^-(g;[v,w])
 =\int_v^w(-g')_+\le\frac{w-v+g(v)-g(w)}2.
\end{equation}

The following potential construction will be used twice. Suppose that
\(H(z)=cz+d\), where \(c\ge1/2\), majorizes the profile at all left
endpoints being considered. Define
\[
 \Lambda(x)=H(2x-1)-g(x).
\]
Then \(\Lambda\) is nondecreasing, and almost everywhere
\begin{equation}\label{minimum:potential}
 (-g')_+\le\frac{2c-g'}{2c+1}=\frac{\Lambda'}{2c+1}.
\end{equation}
Indeed, for \(g'\ge0\) the inequality is immediate, while for
\(g'\le0\) it reduces to \(g'\ge-1\). Through the region
\(\Lambda>0\), use maximal admissible downward edges, shortening the
last edge to the lower boundary of that region. The cost is at most
\(\Lambda(n_0)_+/(2c+1)\), by \eqref{minimum:variation} and
\eqref{minimum:potential}. If \(\Lambda(n)\le0\), then
\[
 g(2n-1)\le H(2n-1)\le g(n),
\]
so moving to a leftmost minimum on \([2n-1,n]\) gives a zero-cost
edge. Each such endpoint is strictly smaller unless the desired terminal
point has already been reached.

In particular, suppose that \(t\) is a minimum of \(g\) on
\([t,1]\), with \(g(t)=m\), and put \(e=g(1)\). The cone
\(H(z)=m+z-t\) majorizes \(g\) on \([t,1]\). Its potential is
nonpositive at \((1+t)/2\). The preceding construction, with \(c=1\),
followed by a free jump to \(t\) as soon as the current endpoint is at
most \((1+t)/2\), gives a chain to \(t\), whenever \(n_0\ge t\),
of cost at most
\begin{equation}\label{minimum:cone}
 \frac{1-t+m-e}{3}.
\end{equation}
This quantity is nonnegative by the Lipschitz inequality. No linear upper
barrier is needed for this cone construction. Zero-cost steps before its
last jump stay in \([t,1]\), since their left endpoints are at least
\(t\).

Here and below, the construction can first be performed for a piecewise
affine profile. A chosen interval minimum is either a breakpoint or the
left endpoint of the interval, except for constant pieces, where the
leftmost convention again chooses a breakpoint or the left endpoint.
There are consequently finitely many steps: steps that do not encounter
a new breakpoint double the complementary depth \(1-n\). The general
continuous case follows by the compactness argument at the end of the
proof. This avoids assuming termination of an arbitrary infinite sequence
of successive minima.

\subsection{The expanded free interval}
For \(n\le q\), a free continuation to zero is always possible. If
\(1/2<n\le q\), then
\[
 g(2n-1)\le b(2n-1)\le an\le g(n).
\]
Choose a leftmost interval minimum to obtain a zero-cost edge. When
\(n\le1/2\), jump directly to zero. This jump is free since
\(g(0)=0\) and \(g\ge0\).

Choose a minimum point \(t\) of \(g\) on \([q,1]\), and write
\[
 m=g(t),\qquad e=g(1).
\]
We shall retain the following relations, rather than immediately replacing
the minimum by the lower barrier:
\begin{equation}\label{minimum:actual-relations}
 t\ge q,\qquad at\le m\le bq,\qquad m\le e,\qquad e\ge a.
\end{equation}
Define
\[
 Q_m=\frac{1+m/b}{2}.
\]
Then \(q\le Q_m\le h\). For \(q<n\le Q_m\),
\[
 g(2n-1)\le b(2n-1)\le m\le g(n).
\]
Thus the free interval actually extends to \(Q_m\), by the same
leftmost-minimum construction and the already proved free continuation
below \(q\).

On \(z\ge m/b\), the affine function
\[
 H_m(z)=\frac z2+\left(b-\frac12\right)\frac mb
\]
majorizes \(bz\). Its potential at \(Q_m\) is
\(m-g(Q_m)\le0\). Apply the potential construction with \(c=1/2\)
down to the expanded free interval, then continue for free. This gives
an always available chain of cost at most
\begin{equation}\label{minimum:floor}
 L(m,e)=\frac{1-2e}{4}-\lambda m
 \le L_0-\lambda m.
\end{equation}
The first expression is nonnegative: it is
\((H_m(1)-g(1))/2\). If the initial point is already in the free
interval, its zero cost satisfies the same estimate.

We will compare this chain with chains whose costs increase with \(m\).
For any \(\sigma>0\), \(\lambda\ge0\), and real \(A\), the elementary
intersection inequality is
\begin{equation}\label{minimum:intersection}
 \min\{L_0-\lambda m,A+\sigma m\}
 \le\frac{\sigma L_0+\lambda A}{\sigma+\lambda}
 \quad(m\in\mathbb R).
\end{equation}
To check it, split at \(m=(L_0-A)/(\sigma+\lambda)\). To the left,
use the increasing line; to the right, use the decreasing line. The
argument includes \(\lambda=0\).

\subsection{An early minimum}
Suppose first that \(t\le h\). In the parameter range
\eqref{minimum:range}, \(q\le2/3\), so \(h\le5/6<7/8\le n_0\).
The cone construction \eqref{minimum:cone} can therefore reach \(t\).
Choose a minimum point \(v\) on \([2t-1,q]\). This interval is
nonempty because \(t\le h\). The edge \(t\to v\) is admissible
and has cost
\[
 (g(v)-m)_+\le\bigl(b(2t-1)-m\bigr)_+.
\]
Indeed, \(g\ge m\) on \([q,t]\), and \(g(v)\) is the minimum
on \([v,q]\). Below \(q\) we continue for free. This yields cost
\begin{equation}\label{minimum:early-route}
 \frac{1-t+m-e}{3}+\bigl(b(2t-1)-m\bigr)_+.
\end{equation}

If the positive part vanishes, \(m\le e\) implies that this cost
is at most \((1-t)/3\le E_1\). Otherwise it is
\[
 A=\frac{1-e-3b+(6b-1)t-2m}{3}.
\]
Since \(e\ge a\), \(m\ge at\), and \(t\le h\),
\begin{equation}\label{minimum:earlyE2}
 A\le\frac{1-a-3b+(6b-1-2a)h}{3}=E_2.
\end{equation}
Alternatively \(t\le m/a\) and \(6b-1>0\) give
\[
 A\le A_E+\sigma_E m.
\]
Comparison with \eqref{minimum:floor}, using
\eqref{minimum:intersection}, bounds the cheaper chain by \(J_E\).
Consequently the early case has cost at most
\begin{equation}\label{minimum:early-bound}
 \max\{E_1,\min(E_2,J_E)\}.
\end{equation}

\subsection{A late minimum and reflection}
Now suppose \(t>h\). Choose a minimum point \(u\) on \([q,h]\),
and put
\[
 w=g(u),\qquad z=b(2u-1).
\]
Here \(q\le u\le h\), \(w\ge m\), and \(2u-1\le q\).
We distinguish whether the terminal left interval is certainly free.

First suppose \(w\le z\). Define a nondecreasing function
\[
 r(x)=\begin{cases}
 0,&x\le u,\\
 w-\min_{[u,x]}g,&x\ge u,
 \end{cases}
 \qquad f(x)=g(x)+r(x).
\]
The reflected function \(f\) is 1-Lipschitz. Here is a direct check.
For \(u\le x<y\), if the running minimum does not decrease on
\([x,y]\), then \(r(x)=r(y)\), so the assertion follows for \(g\).
Otherwise a new minimum is attained at some \(v\in[x,y]\), and
\[
 f(y)=g(y)+w-g(v),\qquad f(x)\ge w.
\]
It follows that \(f(y)-f(x)\le y-v\le y-x\). The lower bound follows
from \(r(y)\ge r(x)\) and \(g(y)-g(x)\ge-(y-x)\). Gluing at
\(u\) handles intervals crossing \(u\).

Moreover, \(f\ge w=f(u)\) on \([q,1]\) and
\(f(1)=e+w-m\). Apply the cone construction to \(f\), reaching
\(u<n_0\) with cost at most
\[
 \frac{1-u-e+m}{3}.
\]
The function is unchanged to the left of \(u\). A terminal jump
to a minimum on \([2u-1,q]\) has cost at most \((z-w)_+\), after
which the original free continuation below \(q\) applies. The
reflection need not satisfy the original linear upper barrier to the
right of \(u\): the cone uses only its Lipschitz property and its
minimum, while the terminal and free portions use its unchanged left part.

For every edge \(n\to v\), monotonicity of \(r\) gives
\[
 c_g(v,n)\le c_f(v,n)+r(n)-r(v).
\]
The extra terms telescope over the entire chain, to at most
\(r(n_0)\le w-m\). We have therefore constructed a chain for \(g\)
of cost at most
\begin{equation}\label{minimum:reflection}
 \frac{1-u-e-2m}{3}+\max(w,z).
\end{equation}
In the present case \(w\le z\), this is
\[
 \frac{1-e-3b+(6b-1)u-2m}{3}\le E_2,
\]
using \(e\ge a\), \(u\le h\), \(m\ge at\ge ah\), and
\(6b-1\ge0\). The floor chain also costs at most
\(L_0-\lambda ah\). At \(m=ah\), the increasing line
\(A_E+\sigma_E m\) has value \(E_2\). Consequently
\eqref{minimum:intersection} implies
\[
 \min\{E_2,L_0-\lambda ah\}\le J_E.
\]
The late case with \(w\le z\) is bounded by \(\min(E_2,J_E)\).

\subsection{A free terminal interval: length and variation}
It remains to consider \(t>h\) and \(w\ge z\). If \(t\le n_0\),
use the cone to reach \(t\), then maximal admissible edges to reach
\(h\). The edge \(h\to u\) is admissible since \(u\ge q=2h-1\),
and it is free since \(u\) minimizes \(g\) on \([q,h]\). The next
edge to a minimum on \([2u-1,q]\) is also free: its cost is at most
\((z-w)_+=0\). Finish with the free continuation below \(q\).

The middle chain from \(t\) to \(h\) has cost at most \(t-h\).
Combining with the cone yields
\[
 \frac{1-t+m-e}{3}+t-h
 =\frac{1-e-3h+2t+m}{3}
 \le S_F+\sigma_F m,
\]
where the last step uses \(e\ge a\) and \(t\le m/a\).
Comparison with the floor chain proves the bound \(J_F\).

For the same middle chain, \eqref{minimum:variation} gives the sharper
alternative
\[
 \operatorname{Var}^-(g;[h,t])
 \le\frac{t-h+g(h)-m}{2}
 \le\frac{t-h+bh-m}{2}.
\]
Its total cost with the cone is therefore at most
\[
 \frac{1-e}{3}-\frac{(1-b)h}{2}+\frac{t-m}{6}
 \le S_V+\sigma_Vm.
\]
Comparison with the floor chain proves the bound \(J_V\).

If \(t>n_0\), omit the cone and start the middle chain at \(n_0\).
Its length is at most \(t-h\) and its negative variation is at most
\(\operatorname{Var}^-(g;[h,t])\). The omitted cone cost
\((1-t+m-e)/3\) is nonnegative. Both displayed upper bounds thus
remain valid. In all cases \(n_0>h\), as checked in the early case.
The late case with \(w\ge z\) is bounded by \(\min(J_F,J_V)\).

\subsection{Removing a redundant term and controlling the chain}
The cases proved so far give cost at most
\[
 \max\{E_1,\min(E_2,J_E),\min(J_F,J_V)\}.
\]
In fact \(E_1\) is redundant. First,
\[
 \frac{E_2}{E_1}=\frac{1+6b-4a}{2}\ge1.
\]
Second, the numerator of \(L_0-\lambda ah-E_1\), after multiplication
by its positive denominator \(24b(2b-a)\), is
\[
 (4-6a)b^2+(6a^2-7a)b+3a^2.
\]
Writing \(b=2a+x\), \(x\ge0\), this is
\[
 a^2(5-12a)+9a(1-2a)x+(4-6a)x^2\ge0.
\]
At \(m=ah\), both lines used to define \(J_E\) consequently have
value at least \(E_1\). Their intersection value is at least their
minimum at any specified point, by \eqref{minimum:intersection}.
Thus \(J_E\ge E_1\), proving precisely the cost
\eqref{minimum:cost}.

Each positive-potential or middle portion consists of maximal edges
\(n\to2n-1\), apart from a bounded number of terminal shortened edges.
Such a maximal edge doubles \(1-n\). Consecutive zero-cost edges
can be merged whenever their union is admissible; the union still has
zero cost because its lower endpoint is a minimum on both constituent
intervals. After all possible adjacent merges, any two consecutive retained
edges, with endpoints \(n_i>n_{i+1}>n_{i+2}\), satisfy
\[
 2n_i-n_{i+2}>1,\quad\text{hence}\quad
 1-n_{i+2}>2(1-n_i).
\]
Every construction above has only a bounded number of such portions and
special terminal edges. It follows that its number of edges is at most
\(C(1+\log(1/(1-n_0)))\).

To justify the continuous assertion fully, interpolate the given profile
on finer and finer partitions containing \(0,n_0,1\). Linear
interpolation preserves both linear barriers and the Lipschitz constant.
Run the finite piecewise affine constructions proved above. Their costs
obey the same bound, and their lengths have the uniform bound just
established. Pad endpoint lists by repeated endpoints and pass to a
subsequence. The profiles converge uniformly. Minima on intervals vary
continuously under uniform convergence and convergence of their endpoints.
Admissibility is closed. The limiting chain therefore has the asserted
cost and length; removing repeated endpoints gives a strictly decreasing
chain.

For a prescribed grid of mesh \(T\), round every endpoint down. If
\(m<n\), \(m\ge2n-N\), and \(N/T\) is integral, then
\[
 \left\lfloor\frac mT\right\rfloor
 \ge \left\lfloor\frac{2n-N}{T}\right\rfloor
 \ge2\left\lfloor\frac nT\right\rfloor-\frac NT.
\]
Thus admissibility is preserved. An endpoint moves by at most \(T\),
and the minimum of a 1-Lipschitz function changes by at most \(T\)
under this change of interval. The increase in cost is at most \(2T\)
per edge. Remove repetitions. Finally, clamping between \(ax\) and
\(bx\) preserves the Lipschitz constant and changes a profile with
approximate barriers uniformly by at most \(\epsilon N\). Each edge
cost changes by at most twice this amount. These arguments prove all
stability and length claims in Lemma~\ref{minimum:profile}.

\subsection{The exact cutoff is a unique root}
\begin{proposition}\label{minimum:root}
For every \(a\in[1/8,1/6]\), there is a unique
\(b_*(a)\in[2a,1/2]\) such that \(V(a,b_*(a))=a\).
The function \(b_*\) is continuous,
\[
 b_*(1/8)=1/4,\qquad b_*(1/6)=1/2,
\]
and \(2a<b_*(a)<1/2\) in the interior. For parameters
\eqref{minimum:range},
\[
 V(a,b)<a\quad\Longleftrightarrow\quad b<b_*(a).
\]
\end{proposition}

\begin{proof}
We first show that \(V(a,b)\) is strictly increasing in \(b\), for
fixed \(a\). Both factors \((b-a)/(2b-a)\) and \(1+6b-4a\)
are positive and strictly increasing, so \(E_2\) is strictly increasing.
Writing \(D_E=24b^2-14ab+3a-4b>0\), differentiation gives
\[
 \frac{\partial J_E}{\partial b}
 =\frac{a(1-a)(60b^2-12b+1-4a)}{D_E^2}>0.
\]
For the last sign, the quadratic in \(b\) is increasing for
\(b\ge2a\ge1/4\); at \(b=2a\) it is
\(240a^2-28a+1\ge5/4\), since its derivative in \(a\) is
positive for \(a\ge1/8\).

For \(J_F,J_V\), the corresponding \(\sigma\) is independent of
\(b\), \(\lambda'(b)<0\), and \(S_F,S_V\) are strictly increasing
in \(b\). Indeed,
\[
 q'(b)=-\frac a{(2b-a)^2}<0,\qquad h'=q'/2<0,
\]
so \(S_F'=-h'>0\), and
\(S_V'=h/2-(1-b)h'/2>0\). Both \(S_F\) and \(S_V\) are less
than \(L_0\). The former is negative; for the latter,
\[
 S_V-L_0=\frac{1+2a}{12}-\frac{(1-b)h}{2}
 \le\frac19-\frac3{16}<0,
\]
using \(h\ge3/4\). Thus the derivative of either intersection value
\(J=(\sigma L_0+\lambda S)/(\sigma+\lambda)\) is
\[
 J'=\frac{\sigma\lambda'(S-L_0)+\lambda(\sigma+\lambda)S'}
              {(\sigma+\lambda)^2}>0.
\]
The minimum or maximum of finitely many strictly increasing functions is
strictly increasing: compare all their values at two ordered arguments.
This proves the strict monotonicity and continuity of \(V\).

At \(b=2a\), direct substitution gives
\[
 E_2=\frac{1+8a}{18}<a,\qquad
 J_F=\frac{5-2a}{2(19-4a)},\qquad
 J_V=\frac{3-2a-16a^2}{4(7-16a)}.
\]
In the stated range \(J_V\le a\), with equality only at \(a=1/8\):
the numerator of \(J_V-a\) is
\(3(1-8a)(1-2a)\). At \(a=1/8\), \(J_F>a\), so
\(V(1/8,1/4)=1/8\); in the interior \(V(a,2a)<a\).
At \(b=1/2\),
\[
 J_E=J_F=J_V=L_0,\qquad E_2=\frac{1-2a}{3}>L_0,
\]
and hence \(V(a,1/2)=(1-2a)/4\ge a\), with equality only at
\(a=1/6\). The intermediate value theorem and strict monotonicity
give the unique root and the claimed endpoint and inequality statements.

For continuity, let \(a_j\to a\) and take any convergent subsequence
of \(b_*(a_j)\), with limit \(b\). The parameter constraints pass
to the limit, and joint continuity gives \(V(a,b)=a\). Uniqueness
implies \(b=b_*(a)\). Compactness then proves continuity of the full
sequence, including at the endpoints.
\end{proof}

\subsection{Comparison, an exact example, and numerical values}
The new bound improves the previous low-slope cost
\[
 U(a,b)=\frac{(b-a)(1-a)}{2(2b-a)}
       =\frac{(1-a)(1-q)}2.
\]
Here is an algebraic check that uses no numerical optimization. At
\(m=aq\), the decreasing floor line has value \(U\); the increasing
early line has value \(2U/3\), while the increasing free-terminal line
has value
\[
 S_F+\sigma_F aq=-\frac{(1+2a)(1-q)}6<0.
\]
Both intersections therefore occur to the right of \(aq\). Their
values satisfy \(J_E,J_F\le U\), strictly when \(b<1/2\) because
then \(\lambda>0\). Consequently
\begin{equation}\label{minimum:old-comparison}
 V(a,b)\le U(a,b),
 \qquad V(a,b)<U(a,b)\quad\hbox{if }b<1/2.
\end{equation}

At \(a=3/20\), \(b=7/20\), exact arithmetic gives
\[
 E_2=\frac5{33},\qquad J_E=\frac{187}{1255},\qquad
 J_F=\frac{1000}{6919},\qquad J_V=\frac{937}{6424}.
\]
It follows by rational comparison that
\[
 V(3/20,7/20)=\frac{187}{1255},\qquad
 \frac3{20}-V(3/20,7/20)=\frac1{1004}>0.
\]
This is the exact profile margin used for the example in the main theorem.
The cutoff at this dimension also has a simple exact expression:
\begin{equation}\label{minimum:exact-cutoff}
 b_*(3/20)=\frac{7+\sqrt{13}}{30},\qquad
 B_{\rm M}(23/20)=\frac{37+\sqrt{13}}{30}.
\end{equation}
To verify the root identification, put \(b=(7+\sqrt{13})/30\).
Substitution gives \(J_E=3/20\) and
\(E_2-3/20=(2\sqrt{13}-7)/60>0\). Also \(b<9/25\), and
strict monotonicity together with the exact calculation
\[
 J_F(3/20,9/25)-3/20=-\frac{136}{51015}<0
\]
shows that \(\min(J_F,J_V)<3/20\). Thus
\(V(3/20,b)=3/20\), and Proposition~\ref{minimum:root} identifies
the claimed unique root.

For convenient evaluation, the level equations for \(E_2,J_E,J_F,J_V\)
to equal \(a\), after clearing positive denominators, are respectively
\begin{align*}
 0={}&6b^2+(1-22a)b+10a^2-a,\\
 0={}&(6-30a)b^2+(20a^2-a-1)b+a-4a^2,\\
 0={}&(8a^2-34a+8)b^2+(-4a^3+4a^2-9a)b+a^2+8a^3,\\
 0={}&-18ab^3+(4-17a+94a^2)b^2
            -(3a+19a^2+44a^3)b+16a^3-a^2.
\end{align*}
One can alternatively evaluate \eqref{minimum:cost} directly and use
monotone bisection. The following decimals are illustrative evaluations of
the unique root; they play no role in the proof or in the definition of
\(b_*\).
\begin{center}
\begin{tabular}{ccc}
\toprule
\(a\)&\(b_*(a)\), approximately&Expression attaining the root\\
\midrule
\(0.125\)&\(0.250000000\)&\(J_V\)\\
\(0.130\)&\(0.269213180\)&\(J_V\)\\
\(0.135\)&\(0.289846951\)&\(J_V\)\\
\(0.140\)&\(0.312168690\)&\(J_V\)\\
\(0.142\)&\(0.322586562\)&\(J_F\)\\
\(0.145\)&\(0.332270658\)&\(E_2\)\\
\(0.150\)&\(0.353518376\)&\(J_E\)\\
\(0.155\)&\(0.386846972\)&\(J_E\)\\
\(0.160\)&\(0.427797338\)&\(J_E\)\\
\(0.165\)&\(0.479534898\)&\(J_E\)\\
\(1/6\)&\(0.500000000\)&\(J_E=J_F=J_V\)\\
\bottomrule
\end{tabular}
\end{center}
Numerically, the active expressions change near
\(a=0.140503259\), \(0.142267709\), and \(0.147793654\), in the
order \(J_V,J_F,E_2,J_E\). These approximate transition locations
are not assertions of an additional piecewise algebraic theorem; the
proved, continuous, unique-root definition already covers the full interval.

\begin{remark}[Distinction from the usual total-drop orientation]
Keleti--Shmerkin's total-drop framework \cite[Definition 5.1 and
Proposition 5.2]{ks} also combines interval minima with dyadically related
endpoints. Its admissible partitions have neighboring scales in ratio at
most two, and its summands subtract an interval minimum from the value at
the smaller scale. Reflecting our coordinate, \(x=1-n\), turns our
admissibility into \(x_{i+1}\le2x_i\), but turns our cost into
\(f(x_{i+1})-\min_{[x_i,x_{i+1}]}f\), which uses the larger endpoint.
The affine barriers are also reflected. Thus that proposition cannot be
substituted with the same parameters for the lemma proved here. Our proof
establishes the required orientation directly; it does not claim an
optimal profile estimate in the present low-slope range.
\end{remark}
